% Part: first-order-logic
% Chapter: syntax-and-semantics
% Section: Semantic Notions

\documentclass[../../../include/open-logic-section]{subfiles}

\begin{document}

\olfileid[hi]{fol}{syn}{sem}
\olsection{अर्थगत धारणाएँ}

\begin{explain}
प्रथम-क्रम भाषाओं के लिए !!{structure} की परिभाषा मिलने पर हम वाक्यों के
कुछ मूल अर्थगत गुण और उनके बीच के संबंध परिभाषित कर सकते हैं। इनमें सबसे
सरल किसी वाक्य की \emph{वैधता} की धारणा है। कोई वाक्य वैध है यदि वह प्रत्येक
!!a{structure} में संतुष्ट है। वैध वाक्य इस बात से निरपेक्ष होकर संतुष्ट
होते हैं कि उनमें आए अतार्किक चिह्नों का निर्वचन किस प्रकार किया गया है।
इसलिए वैध वाक्यों को \emph{तार्किक सत्य} भी कहते हैं---वे किसी भी
!!a{structure} में सत्य, अर्थात् संतुष्ट, होते हैं; अतः उनकी सत्यता केवल
उनमें आए तार्किक चिह्नों और उनकी वाक्यविन्यासगत !!{structure} पर निर्भर करती
है, अतार्किक चिह्नों अथवा उनके निर्वचन पर नहीं।
\end{explain}

\begin{defn}[वैधता]
कोई वाक्य~$!A$ \emph{वैध} है, $\Entails !A$, तभी और केवल तभी जब
$\Sat{M}{!A}$, प्रत्येक !!a{structure}~$\Struct M$ के लिए।
\end{defn}

\begin{defn}[तात्पर्य]
वाक्यों का कोई समुच्चय~$\Gamma$, किसी वाक्य~$!A$ से \emph{तात्पर्य रखता
है}, $\Gamma
\Entails !A$, तभी और केवल तभी जब प्रत्येक ऐसी
!!a{structure}~$\Struct M$ के लिए, जिसके लिए $\Sat{M}{\Gamma}$,
$\Sat{M}{!A}$ भी।
\end{defn}


\begin{defn}[संतुष्टनीयता]
वाक्यों का कोई समुच्चय~$\Gamma$ \emph{संतुष्टनीय} है यदि $\Sat{M}{\Gamma}$
किसी !!a{structure}~$\Struct M$ के लिए। यदि~$\Gamma$
संतुष्टनीय नहीं है, तो उसे \emph{असंतुष्टनीय} कहते हैं।
\end{defn}

\begin{prop}
कोई वाक्य~$!A$ वैध है तभी और केवल तभी जब $\Gamma \Entails !A$,
वाक्यों के प्रत्येक समुच्चय~$\Gamma$ के लिए।
\end{prop}

\begin{proof}
अग्र दिशा के लिए मान लें कि~$!A$ वैध है, और~$\Gamma$ वाक्यों का कोई
समुच्चय है। ऐसी !!a{structure}~$\Struct M$ लें कि $\Sat{M}{\Gamma}$।
चूँकि~$!A$ वैध है, $\Sat{M}{!A}$; अतः $\Gamma
\Entails !A$।

विलोम दिशा के प्रतिपक्ष के लिए मान लें कि~$!A$ अवैध है; अतः कोई
!!a{structure}~$\Struct M$ ऐसी है कि $\Sat/{M}{!A}$। जब $\Gamma
= \{ \ltrue \}$, तब~$\ltrue$ के वैध होने से $\Sat{M}{\Gamma}$। अतः कोई
!!a{structure}~$\Struct M$ ऐसी है कि $\Sat{M}{\Gamma}$ किंतु
$\Sat/{M}{!A}$; इसलिए~$\Gamma$,~$!A$ से तात्पर्य नहीं रखता।
\end{proof}

\begin{prop}
\ollabel{prop:entails-unsat}
$\Gamma \Entails !A$ तभी और केवल तभी जब $\Gamma \cup \{\lnot !A\}$
असंतुष्टनीय है।
\end{prop}

\begin{proof}
अग्र दिशा के लिए मान लें कि~$\Gamma \Entails !A$, और इसके विपरीत मान लें
कि कोई !!a{structure}~$\Struct M$ ऐसी है कि
$\Sat{M}{\Gamma
  \cup \{ \lnot !A \}}$। चूँकि $\Sat{M}{\Gamma}$ और $\Gamma \Entails
!A$, इसलिए $\Sat{M}{!A}$। साथ ही, चूँकि
$\Sat{M}{\Gamma\cup \{ \lnot !A \}}$, इसलिए $\Sat{M}{\lnot !A}$;
अतः हमारे पास $\Sat{M}{!A}$ और $\Sat/{M}{!A}$ दोनों हैं, जो विरोधाभास है।
इसलिए ऐसी कोई !!{structure}~$\Struct M$ नहीं हो सकती, और
$\Gamma \cup \{ \lnot !A \}$ असंतुष्टनीय है।

विलोम दिशा के लिए मान लें कि $\Gamma \cup \{ \lnot !A \}$
असंतुष्टनीय है। अतः प्रत्येक !!a{structure}~$\Struct M$ के लिए या तो
$\Sat/{M}{\Gamma}$ अथवा $\Sat{M}{!A}$। इसलिए प्रत्येक ऐसी !!{structure}
$\Struct M$ के लिए, जिसके लिए $\Sat{M}{\Gamma}$, $\Sat{M}{!A}$; अतः
$\Gamma \Entails
!A$।
\end{proof}

\begin{prob}
\begin{enumerate}
\item दिखाइए कि $\Gamma \Entails \bot$ तभी और केवल तभी जब~$\Gamma$
  असंतुष्टनीय है।
\item दिखाइए कि $\Gamma \cup \{!A\} \Entails \bot$ तभी और केवल तभी जब
  $\Gamma \Entails \lnot !A$।
\item मान लें कि~$c$,~$!A$ अथवा~$\Gamma$ में नहीं आता। दिखाइए कि
  $\Gamma \Entails \lforall[x][!A]$ तभी और केवल तभी जब $\Gamma \Entails
  \Subst{!A}{c}{x}$।
\end{enumerate}
\end{prob}

\begin{prop}
यदि $\Gamma \subseteq \Gamma'$ और $\Gamma \Entails !A$, तो
$\Gamma'
\Entails !A$।
\end{prop}

\begin{proof}
मान लें कि $\Gamma \subseteq \Gamma'$ और $\Gamma \Entails !A$। ऐसी
संरचना~$\Struct M$ लें कि $\Sat{M}{\Gamma'}$; तब $\Sat{M}{\Gamma}$, और
चूँकि $\Gamma \Entails !A$, हमें $\Sat{M}{!A}$ मिलता है। अतः जब भी
$\Sat{M}{\Gamma'}$, तब $\Sat{M}{!A}$; इसलिए $\Gamma' \Entails !A$।
\end{proof}


\begin{thm}[अर्थगत निगमन प्रमेय]
\ollabel{thm:sem-deduction}
$\Gamma \cup \{!A\} \Entails !B$ तभी और केवल तभी जब
$\Gamma \Entails !A \lif !B$।
\end{thm}

\begin{proof}
अग्र दिशा के लिए मान लें कि $\Gamma \cup \{ !A \} \Entails !B$, और ऐसी
!!a{structure}~$\Struct M$ लें कि $\Sat{M}{\Gamma}$। यदि $\Sat{M}{!A}$,
  तो $\Sat{M}{\Gamma \cup \{ !A \} }$; अतः चूँकि $\Gamma
\cup \{ !A \}$ से~$!B$ का तात्पर्य है, हमें $\Sat{M}{!B}$ मिलता है। इसलिए
$\Sat{M}{!A \lif !B}$, और अतः $\Gamma \Entails !A \lif !B$।

विलोम दिशा के लिए मान लें कि $\Gamma \Entails !A \lif !B$, और ऐसी
!!a{structure}~$\Struct M$ लें कि $\Sat{M}{\Gamma \cup \{ !A
  \}}$। तब
$\Sat{M}{\Gamma}$, अतः $\Sat{M}{!A \lif !B}$; और चूँकि $\Sat{M}{!A}$,
इसलिए $\Sat{M}{!B}$। अतः जब भी $\Sat{M}{\Gamma \cup \{
  !A \} }$, तब
$\Sat{M}{!B}$; इसलिए $\Gamma \cup \{ !A \} \Entails !B$।
\end{proof}

\begin{prop}\ollabel{prop:quant-terms}
  मान लें कि~$\Struct{M}$ कोई !!a{structure},~$!A(x)$ एक मुक्त चर~$x$ वाला
  कोई !!a{formula}, और~$t$ कोई बंद पद है। तब:
  \begin{tagenumerate}{prvEx,prvAll}
    \tagitem{prvEx}{$!A(t) \Entails \lexists[x][!A(x)]$}{}
    \tagitem{prvAll}{$\lforall[x][!A(x)] \Entails !A(t)$}{}
  \end{tagenumerate}
\end{prop}

\begin{proof}
  \begin{tagenumerate}{prvEx,prvAll}
    \tagitem{prvEx}{%
      \iftag{probEx}{अभ्यास।}{मान लें कि $\Sat{M}{!A(t)}$। ऐसा चर-नियतन
        ~$s$ लें कि $s(x) = \Value{t}{M}$। तब $\Sat{M}{!A(t)}[s]$, क्योंकि
        $!A(t)$ कोई !!a{sentence} है।
        \olref[ext]{prop:ext-formulas} से $\Sat{M}{!A(x)}[s]$।
        \olref[ass]{prop:sat-quant} से
        $\Sat{M}{\lexists[x][!A(x)]}$।}}{}
    \tagitem{prvAll}{%
      \iftag{probAll}{अभ्यास।}{मान लें कि
        $\Sat{M}{\lforall[x][!A(x)]}$। ऐसा चर-नियतन~$s$ लें कि
        $s(x) = \Value{t}{M}$। \olref[ass]{prop:sat-quant} से
        $\Sat{M}{!A(x)}[s]$। \olref[ext]{prop:ext-formulas} से
        $\Sat{M}{!A(t)}[s]$। \olref[ass]{prop:sentence-sat-true} से
        $\Sat{M}{!A(t)}$, क्योंकि $!A(t)$ कोई !!a{sentence} है।}}{}
  \end{tagenumerate}
\end{proof}

\begin{probtag}{probEx,probAll}
  \olref[fol][syn][sem]{prop:quant-terms} का प्रमाण पूरा कीजिए।
\end{probtag}

\end{document}
